\documentclass[11pt]{article}
\usepackage[margin=1in]{geometry}
\usepackage[T1]{fontenc}
\usepackage{lmodern}
\usepackage[hidelinks]{hyperref}
\usepackage{booktabs}
\usepackage{microtype}

\title{Concrete Conflict-Analysis Calculus and Guide}
\author{Tactix Development Project}
\date{August 8, 2026}

\begin{document}
\maketitle

\section*{Executive summary}
The M3 ConflictAnalysis entry in the interactive proof-rule guide previously
compressed learning, backjumping, assertion, and proof logging into informal
placeholders.  The authoritative calculus and the derived HTML guide now give
the exact first-UIP construction, jump-level definition, atomic trail
transition, and ordered certificate delta.  This is a documentation-only
clarification; solver semantics and performance are unchanged.

\section{Findings and results}
The old display used ``derive learned clause'', \texttt{backtrack(j)}, and
``optional deletion'' without defining the learned clause, the level $j$, or
which deletions occur.  The revised rule starts from the all-false conflict
clause, resolves the latest conflict-level trail assignments with their reason
clauses until one conflict-level literal remains, and identifies that literal
as the negated first UIP.  It defines the jump level as the maximum level among
the other learned literals (or zero), and states the legal chronological target
interval.

The HTML certificate is now ordered explicitly: OTFS add-before-delete pairs,
the learned RUP addition with its reason replay, and eager-subsumption
deletions.  Root LRAT and DRAT refutations are distinguished.  A new two-level
worked example shows both resolution steps, the LRAT hint order, the computed
jump level, and the asserted literal.  The root-refutation example now shows
its concrete certificate as well.

\section{Evidence}
The statement was checked against \path{src/rules/analyze/analyze.cpp},
\path{src/rules/backjump/backjump.cpp}, and the proof API.  Documentation
citation, section-reference, size, comment-hygiene, and whitespace checks pass.
No solver binary was changed, so no correctness or performance claim is made
from this work.

\section{Risks and decisions}
The guide remains a derived teaching view; \texttt{docs/calculus.md} remains
authoritative.  Only M3 was expanded in this increment.  Several other cards
still use compact quantified summaries such as ``rewrites'' or ``resolvents'';
they should be made concrete individually against their authoritative calculus
and implementation rather than expanded speculatively in one broad rewrite.

\section{Next actions}
Review the M3 presentation in a browser.  If this level of detail is preferred,
apply the same notation pattern to the remaining cards with aggregate proof
deltas, beginning with Substitute, Sweep, Congruence, and BVE.  The solver's
development next step remains the \texttt{SP-COMPACT-1} Kissat-parity row.

\end{document}
